% Options for packages loaded elsewhere
% Options for packages loaded elsewhere
\PassOptionsToPackage{unicode}{hyperref}
\PassOptionsToPackage{hyphens}{url}
\PassOptionsToPackage{dvipsnames,svgnames,x11names}{xcolor}
%
\documentclass[
  letterpaper,
  DIV=11,
  numbers=noendperiod]{scrartcl}
\usepackage{xcolor}
\usepackage{amsmath,amssymb}
\setcounter{secnumdepth}{-\maxdimen} % remove section numbering
\usepackage{iftex}
\ifPDFTeX
  \usepackage[T1]{fontenc}
  \usepackage[utf8]{inputenc}
  \usepackage{textcomp} % provide euro and other symbols
\else % if luatex or xetex
  \usepackage{unicode-math} % this also loads fontspec
  \defaultfontfeatures{Scale=MatchLowercase}
  \defaultfontfeatures[\rmfamily]{Ligatures=TeX,Scale=1}
\fi
\usepackage{lmodern}
\ifPDFTeX\else
  % xetex/luatex font selection
\fi
% Use upquote if available, for straight quotes in verbatim environments
\IfFileExists{upquote.sty}{\usepackage{upquote}}{}
\IfFileExists{microtype.sty}{% use microtype if available
  \usepackage[]{microtype}
  \UseMicrotypeSet[protrusion]{basicmath} % disable protrusion for tt fonts
}{}
\makeatletter
\@ifundefined{KOMAClassName}{% if non-KOMA class
  \IfFileExists{parskip.sty}{%
    \usepackage{parskip}
  }{% else
    \setlength{\parindent}{0pt}
    \setlength{\parskip}{6pt plus 2pt minus 1pt}}
}{% if KOMA class
  \KOMAoptions{parskip=half}}
\makeatother
% Make \paragraph and \subparagraph free-standing
\makeatletter
\ifx\paragraph\undefined\else
  \let\oldparagraph\paragraph
  \renewcommand{\paragraph}{
    \@ifstar
      \xxxParagraphStar
      \xxxParagraphNoStar
  }
  \newcommand{\xxxParagraphStar}[1]{\oldparagraph*{#1}\mbox{}}
  \newcommand{\xxxParagraphNoStar}[1]{\oldparagraph{#1}\mbox{}}
\fi
\ifx\subparagraph\undefined\else
  \let\oldsubparagraph\subparagraph
  \renewcommand{\subparagraph}{
    \@ifstar
      \xxxSubParagraphStar
      \xxxSubParagraphNoStar
  }
  \newcommand{\xxxSubParagraphStar}[1]{\oldsubparagraph*{#1}\mbox{}}
  \newcommand{\xxxSubParagraphNoStar}[1]{\oldsubparagraph{#1}\mbox{}}
\fi
\makeatother


\usepackage{longtable,booktabs,array}
\newcounter{none} % for unnumbered tables
\usepackage{calc} % for calculating minipage widths
% Correct order of tables after \paragraph or \subparagraph
\usepackage{etoolbox}
\makeatletter
\patchcmd\longtable{\par}{\if@noskipsec\mbox{}\fi\par}{}{}
\makeatother
% Allow footnotes in longtable head/foot
\IfFileExists{footnotehyper.sty}{\usepackage{footnotehyper}}{\usepackage{footnote}}
\makesavenoteenv{longtable}
\usepackage{graphicx}
\makeatletter
\newsavebox\pandoc@box
\newcommand*\pandocbounded[1]{% scales image to fit in text height/width
  \sbox\pandoc@box{#1}%
  \Gscale@div\@tempa{\textheight}{\dimexpr\ht\pandoc@box+\dp\pandoc@box\relax}%
  \Gscale@div\@tempb{\linewidth}{\wd\pandoc@box}%
  \ifdim\@tempb\p@<\@tempa\p@\let\@tempa\@tempb\fi% select the smaller of both
  \ifdim\@tempa\p@<\p@\scalebox{\@tempa}{\usebox\pandoc@box}%
  \else\usebox{\pandoc@box}%
  \fi%
}
% Set default figure placement to htbp
\def\fps@figure{htbp}
\makeatother


% definitions for citeproc citations
\NewDocumentCommand\citeproctext{}{}
\NewDocumentCommand\citeproc{mm}{%
  \begingroup\def\citeproctext{#2}\cite{#1}\endgroup}
\makeatletter
 % allow citations to break across lines
 \let\@cite@ofmt\@firstofone
 % avoid brackets around text for \cite:
 \def\@biblabel#1{}
 \def\@cite#1#2{{#1\if@tempswa , #2\fi}}
\makeatother
\newlength{\cslhangindent}
\setlength{\cslhangindent}{1.5em}
\newlength{\csllabelwidth}
\setlength{\csllabelwidth}{3em}
\newenvironment{CSLReferences}[2] % #1 hanging-indent, #2 entry-spacing
 {\begin{list}{}{%
  \setlength{\itemindent}{0pt}
  \setlength{\leftmargin}{0pt}
  \setlength{\parsep}{0pt}
  % turn on hanging indent if param 1 is 1
  \ifodd #1
   \setlength{\leftmargin}{\cslhangindent}
   \setlength{\itemindent}{-1\cslhangindent}
  \fi
  % set entry spacing
  \setlength{\itemsep}{#2\baselineskip}}}
 {\end{list}}
\usepackage{calc}
\newcommand{\CSLBlock}[1]{\hfill\break\parbox[t]{\linewidth}{\strut\ignorespaces#1\strut}}
\newcommand{\CSLLeftMargin}[1]{\parbox[t]{\csllabelwidth}{\strut#1\strut}}
\newcommand{\CSLRightInline}[1]{\parbox[t]{\linewidth - \csllabelwidth}{\strut#1\strut}}
\newcommand{\CSLIndent}[1]{\hspace{\cslhangindent}#1}



\setlength{\emergencystretch}{3em} % prevent overfull lines

\providecommand{\tightlist}{%
  \setlength{\itemsep}{0pt}\setlength{\parskip}{0pt}}



 


\KOMAoption{captions}{tableheading}
\makeatletter
\@ifpackageloaded{caption}{}{\usepackage{caption}}
\AtBeginDocument{%
\ifdefined\contentsname
  \renewcommand*\contentsname{Table of contents}
\else
  \newcommand\contentsname{Table of contents}
\fi
\ifdefined\listfigurename
  \renewcommand*\listfigurename{List of Figures}
\else
  \newcommand\listfigurename{List of Figures}
\fi
\ifdefined\listtablename
  \renewcommand*\listtablename{List of Tables}
\else
  \newcommand\listtablename{List of Tables}
\fi
\ifdefined\figurename
  \renewcommand*\figurename{Figure}
\else
  \newcommand\figurename{Figure}
\fi
\ifdefined\tablename
  \renewcommand*\tablename{Table}
\else
  \newcommand\tablename{Table}
\fi
}
\@ifpackageloaded{float}{}{\usepackage{float}}
\floatstyle{ruled}
\@ifundefined{c@chapter}{\newfloat{codelisting}{h}{lop}}{\newfloat{codelisting}{h}{lop}[chapter]}
\floatname{codelisting}{Listing}
\newcommand*\listoflistings{\listof{codelisting}{List of Listings}}
\makeatother
\makeatletter
\makeatother
\makeatletter
\@ifpackageloaded{caption}{}{\usepackage{caption}}
\@ifpackageloaded{subcaption}{}{\usepackage{subcaption}}
\makeatother
\usepackage{bookmark}
\IfFileExists{xurl.sty}{\usepackage{xurl}}{} % add URL line breaks if available
\urlstyle{same}
% fallback for those not using the hyperref driver hyperxmp:
\makeatletter
\@ifundefined{xmpquote}{\newcommand{\xmpquote}[1]{#1}}{}
\makeatother
\hypersetup{
  pdftitle={Wildfires and Power Outages? How do they interact? Do they interact? Let's find out! (Replace me)},
  colorlinks=true,
  linkcolor={blue},
  filecolor={Maroon},
  citecolor={Blue},
  urlcolor={Blue},
  pdfcreator={LaTeX via pandoc}}


\title{Wildfires and Power Outages? How do they interact? Do they
interact? Let's find out! (Replace me)}
\author{Heather McBrien \and Logan Piepmeier \and Joan Casey}
\date{}
\begin{document}
\maketitle


\subsection{Abstract}\label{abstract}

\emph{Write me.}

\subsection{Introduction}\label{introduction}

\emph{Write me.}

\subsection{Methods}\label{methods}

\subsubsection{Data}\label{data}

\paragraph{Power Outages}\label{power-outages}

We determined when a county experienced power outage using the
Environment for the Analysis of Geo-Located Energy Information (EAGLE-I)
data set(Brelsford et al. 2023), which records the number of customers
per US county experiencing power outages in 15 minute intervals. We
aggregated that data to daily intervals, defining a county to be
affected by power outage on a particular day when at least \textbf{XXX\%
FIX ME ONCE WE DECIDE} or 5,000 customers had been without power for at
least eight hours. I.e. for a power outage from 10:00 PM on day one to
6:00 AM on day two, we considered the county to have experienced an
outage on day two, when the length of the outage first reached eight
hours. For a longer outage from 8:00 AM on day 1 to 9:00 PM on day two,
we considered both days to be affected by outage since day one had
accumulated eight consecutive hours of outage by 4:00 PM, and day two
had experienced eight consecutive hours as soon as it began on midnight.

The total number of customers per county was not available for all
years. When it was missing, we estimated it by using the number of
households in the American Community Survey (ACS) for that year and
county. ACS data for 2025 was not yet available, so we used the 2024
household count for that year. ACS estimates were obtained with the
Tidycensus R Package(Walker and Herman 2026).

Although the EAGLE-I data set covers the United States and was available
for 2014 through 2025, its geographic coverage improved considerably in
2018. We therefore restricted our analysis to 2018 and later.

\paragraph{Wildfire Smoke}\label{wildfire-smoke}

We determined when a county was affected by wilfire smoke using a model
of daily ambient PM\textsubscript{2.5} attributable to wildfire smoke
(WFS)(Childs et al. 2025). These WFS predictions were available
aggregated to county level. The model identifies a day as a smoke day if
a if the NOAA Hazard Mapping System satellite data indicated smoke plume
overhead, regardless of the modeled concentration of
PM\textsubscript{2.5}. We used a modfied definition for a smoke day,
considering a day to be exposed to WFS only if on smoke days with WFS
PM\textsubscript{2.5} concentrations of \textbf{XXX FIX ME} µg/m³ or
above.

This model covered the continental United States for 2008 through 2024.

\paragraph{Wildfire Burn Zones}\label{wildfire-burn-zones}

We determined whether a county was affected by a wildfire disaster using
a data set of wildfire burn zones (WFBZ) for fires resulting in at least
one death, one destroyed building, or a FEMA major disaster declaration
(Wilner et al. 2025). The data set includes the largest extent of these
fires' perimeters, the estimated ignition date, and, when available, the
containment date. For fires without a containment date, we estimated it
to be four days after ignition. Four days was the 25th percentile of all
known fire durations in the data set.

We obtained 100 meter gridded population estimates from the Global Human
Settlement Layer(Pesaresi et al. 2024), and used these to calculate the
number of people living within 10 km of each fire and living in each
county. We considered a county to be affected when at least
\textbf{XXX\% FIX ME ONCE WE DECIDE} of the county population or 5,000
people lived within 10 km of any active disaster fire on that day. This
data set covered the United States from 2000 through 2025.

\paragraph{Combining Data Sets}\label{combining-data-sets}

We then merged the three data sets on county and day, restricting to
days covered by all three input data sets: January 1, 2018 through
December 31, 2024. Each county and day thus had a binary value for:

\begin{itemize}
\tightlist
\item
  power outage
\item
  exposure to wildfire smoke
\item
  proximity to wildfire disaster
\end{itemize}

\subsubsection{Visualization}\label{visualization}

\emph{We made a pretty map!}

\subsection{Results}\label{results}

\emph{Here is our pretty map!}

\subsection{Discussion}\label{discussion}

\emph{Our map is great!}

\subsection{References}\label{references}

\protect\phantomsection\label{refs}
\begin{CSLReferences}{1}{1}
\bibitem[\citeproctext]{ref-Brelsford2023}
Brelsford, Christa, Sarah Tennille, Aaron Myers, et al. 2023. \emph{{The
Environment for Analysis of Geo-Located Energy Information's Recorded
Electricity Outages 2014-2025}}. November.
\url{https://doi.org/10.6084/m9.figshare.24237376.v4}.

\bibitem[\citeproctext]{ref-Childs2025-hu}
Childs, Marissa, Mariana Martins, Andrew Wilson, Minghao Qiu, Sam
Heft-Neal, and Marshall Burke. 2025. {``Growing Wildfire-Derived {PM2.5}
Across the Contiguous {U.S}. And Implications for Air Quality
Regulation.''} In \emph{EarthArXiv}. August.

\bibitem[\citeproctext]{ref-Pesaresi2024-aq}
Pesaresi, Martino, Marcello Schiavina, Panagiotis Politis, et al. 2024.
{``Advances on the Global Human Settlement Layer by Joint Assessment of
Earth Observation and Population Survey Data.''} \emph{Int. J. Digit.
Earth} 17 (1).

\bibitem[\citeproctext]{ref-Walker2026}
Walker, Kyle, and Matt Herman. 2026. \emph{Tidycensus: Load US Census
Boundary and Attribute Data as 'Tidyverse' and 'Sf'-Ready Data Frames}.
\url{https://walker-data.com/tidycensus/}.

\bibitem[\citeproctext]{ref-Wilner2025-kr}
Wilner, Lauren B, Logan Piepmeier, Milo Gordon, et al. 2025. {``Two and
a Half Decades of United States Wildfire Burn Zone Disaster Data,
2000-2025.''} \emph{Sci. Data} 12 (1): 1948.

\end{CSLReferences}




\end{document}
